\section{Formatos de imagen cruda} \label{sec:appendix_D_raw_formats}

Este apéndice resume aspectos relevantes sobre
formatos de imagen cruda desde una perspectiva práctica, centrándose en 
sus características fundamentales, su uso en distintos contextos y las
implicaciones que tienen a nivel de implementación.

\subsection{Generalidades de los formatos crudos}

En términos generales, una imagen cruda almacena directamente las muestras del sensor con el menor
nivel posible de procesado en cámara. A diferencia de formatos finales como
\texttt{.jpg}, el dato crudo prioriza fidelidad radiométrica frente a tamaño o
portabilidad inmediata.

El caso más representativo es \texttt{.raw} "genérico": un volcado binario de
píxeles en el orden de adquisición, sin estructura estandarizada de metadatos.
Por este motivo, para interpretar correctamente el fichero suelen ser necesarios,
además de los datos de píxel, parámetros como:

\begin{itemize}[noitemsep]
    \item ancho y alto de imagen,
    \item profundidad de bits efectiva (por ejemplo, 10, 12, 14 o 16 bits),
    \item patrón Bayer (\texttt{RGGB}, \texttt{GRBG}, \texttt{GBRG}, \texttt{BGGR}),
    \item endianness (\textit{little} o \textit{big endian}),
    \item posible empaquetado de bits por píxel.
\end{itemize}

Estos parámetros pueden ser proporcionados por el sistema de adquisición, inferidos a
partir de convenciones de nomenclatura o almacenados en un archivo auxiliar de
metadatos.

\subsection{Consideraciones prácticas y curiosidades}

En entornos industriales y científicos es habitual que el fichero de imagen se
entregue junto con un archivo auxiliar de metadatos (\texttt{.ini}, \texttt{.cfg},
\texttt{.txt}, \texttt{.json}, \texttt{.xml}, \texttt{.hdr}, etc.). Esta
estrategia simplifica la adquisición y reduce dependencias de formato, aunque
obliga al \textit{software} de carga a mantener coherencia estricta entre ambos
ficheros.

En sistemas embebidos, cuando la cadena de adquisición es fija, es común
asumir dimensiones y configuración constantes por aplicación. Esto permite
maximizar rendimiento y minimizar complejidad de procesamiento, a costa de reducir
flexibilidad ante cambios de sensor o perfil de captura.

Por otro lado, en fotografía digital profesional predominan formatos propietarios
de fabricante (por ejemplo, \texttt{.CR2}, \texttt{.CR3}, \texttt{.NEF},
\texttt{.ARW}, \texttt{.RAF}, \texttt{.ORF}, \texttt{.RW2}), que sí integran
metadatos dentro del propio fichero que contiene los píxeles de la imagen. 
En estos casos, la lectura suele apoyarse en bibliotecas específicas de \textit{decoding} RAW.

\newpage

\subsection{Formatos crudos más comunes}

Las \autoref{tab:appendix_d_raw_formats_common_i} y
\autoref{tab:appendix_d_raw_formats_common_ii} resumen los formatos de imagen
cruda más representativos en la práctica, su contexto típico de uso y la forma en
que se gestionan los metadatos.

\begin{table}[htbp]
    \centering
    \footnotesize
    \renewcommand{\arraystretch}{1.15}
    \setlength{\tabcolsep}{4pt}
    \begin{tabularx}{\textwidth}{|>{\hsize=0.95\hsize\raggedright\arraybackslash}X|>{\hsize=1.00\hsize\raggedright\arraybackslash}X|>{\hsize=0.85\hsize\raggedright\arraybackslash}X|>{\hsize=1.20\hsize\raggedright\arraybackslash}X|}
        \hline
        \rowcolor[gray]{0.90}
        \textbf{Formato / Extensión} & \textbf{Ámbito típico} & \textbf{Metadatos} & \textbf{Notas de uso} \\
        \hline
        \texttt{.raw} (genérico) & Visión industrial, sensores embebidos, captura científica & Externos o implícitos en la configuración del sistema & Máxima simplicidad de almacenamiento; requiere conocer geometría, bits y patrón de color para decodificar. \\
        \hline
        \texttt{.dng} (Digital Negative) & Fotografía y archivado interoperable & Integrados en el propio fichero & Alternativa abierta con amplia compatibilidad; facilita preservación y trazabilidad. \\
        \hline
        \texttt{.ari}, \texttt{.3fr}, \texttt{.fff}, \texttt{.iiq}, \texttt{.x3f} & Cine digital y captura profesional de alta gama & Integrados (propietarios) & Orientados a máxima calidad y rango dinámico; suelen implicar flujos de trabajo especializados. \\
        \hline
        \texttt{.raw} + \textit{sidecar} (\texttt{.ini}, \texttt{.cfg}, \texttt{.json}, \texttt{.xml}, \texttt{.hdr}) & Microscopía, teledetección e instrumentación científica & Externos en fichero auxiliar & Patrón frecuente cuando se prioriza reproducibilidad experimental y control explícito de parámetros. \\
        \hline
    \end{tabularx}
    \caption{Resumen de formatos de imagen cruda y contexto de uso (I)}
    \label{tab:appendix_d_raw_formats_common_i}
\end{table}

\begin{table}[htbp]
    \centering
    \footnotesize
    \renewcommand{\arraystretch}{1.15}
    \setlength{\tabcolsep}{4pt}
    \begin{tabularx}{\textwidth}{|>{\hsize=0.75\hsize\raggedright\arraybackslash}X|>{\hsize=1.35\hsize\raggedright\arraybackslash}X|>{\hsize=0.90\hsize\raggedright\arraybackslash}X|}
        \hline
        \rowcolor[gray]{0.90}
        \textbf{Fabricante / ecosistema} & \textbf{Extensiones RAW más comunes} & \textbf{Observación práctica} \\
        \hline
        Canon & \texttt{.cr2}, \texttt{.cr3}, \texttt{.crw} & Evolución de formatos propietarios con cambios de estructura entre generaciones. \\
        \hline
        Nikon & \texttt{.nef}, \texttt{.nrw} & Ampliamente soportado por librerías de revelado RAW y herramientas de catálogo. \\
        \hline
        Sony & \texttt{.arw}, \texttt{.sr2}, \texttt{.srf} & Presencia habitual en flujos híbridos foto-vídeo de alta resolución. \\
        \hline
        Fujifilm / Olympus / Panasonic / Pentax & \texttt{.raf}, \texttt{.orf}, \texttt{.rw2}, \texttt{.pef} & Formatos consolidados en fotografía, con semánticas internas específicas por marca. \\
        \hline
        Otros fabricantes y nichos & \texttt{.dcr}, \texttt{.kdc}, \texttt{.mos}, \texttt{.rwl}, \texttt{.mef}, \texttt{.mrw}, \texttt{.srw}, etc. & La heterogeneidad de metadatos y compresión refuerza la necesidad de \textit{decoders} especializados. \\
        \hline
    \end{tabularx}
    \caption{Listado de extensiones RAW propietarias más comunes (II)}
    \label{tab:appendix_d_raw_formats_common_ii}
\end{table}

Como referencia operativa dentro de este TFM, los \textit{loaders}
implementados en la arquitectura del sistema (\autoref{subsubsec:7_3_2_loaders}) dan soporte directo a
\texttt{.raw}, \texttt{.arw} y \texttt{.dng}, mientras que
\texttt{.pkl.gz} y \texttt{.jpg} se incluyen como formatos auxiliares para
ensayo, validación y pruebas de integración.
